\chapter{系统实现}

在完成方法设计与正式实验之后，还需要说明本文提出的检索方案如何被实现为可运行、可复现、可演示的工程系统。因此，本章从工程实现角度对后端服务、索引构建、离线评测、前端交互以及系统可运行性进行说明，展示本文工作不仅停留在算法设计层面，而且已经形成较完整的原型系统。

\section{实现环境}

系统后端采用 Python 3.10 开发，模型推理运行于 CUDA 环境，服务框架为 FastAPI，向量数据库使用 Qdrant，OCR 模块采用 RapidOCR 方案。围绕本文所提出的方法体系，系统实现已完成以下几项关键功能建设：
\begin{enumerate}
  \item 支持面向 \texttt{Dataset\_img3} 的全图库重建；
  \item 支持独立建库脚本 \texttt{build\_index.py}，可显式指定 \texttt{CUDA} 设备并按需关闭 OCR；
  \item 支持离线模式对齐评测脚本，用于在不同模式下分别重建图库和查询特征；
  \item 支持前后端一体化检索服务与结果展示界面。
\end{enumerate}

\section{核心模块实现}

当前系统主要由以下模块构成：
\begin{enumerate}
  \item \texttt{feature\_service.py}：负责视觉特征提取、掩码引导聚合、结构描述符构建以及 OCR 辅助所需的基础特征组织；
  \item \texttt{cleaning\_service.py}：负责基于 YOLO 的图纸清洗和保留掩码生成；
  \item \texttt{ocr\_service.py}：负责 OCR 文本抽取、标题栏字段结构化、字段相似度计算；
  \item \texttt{qdrant\_db.py}：负责图库向量写入、候选召回、结构重排序和在线检索结果组织；
  \item \texttt{build\_index.py}：负责 GPU 化独立建库，避免在服务启动阶段直接触发高负载全量建库；
  \item \texttt{evaluate\_visual\_offline.py}：负责视觉主线离线对齐评测；
  \item \texttt{evaluate\_multimodal\_offline.py}：负责 OCR 辅助模式离线评测；
  \item \texttt{visualize\_pair\_similarity.py}：负责相似案例分析与 Grad-CAM 可视化。
\end{enumerate}

各模块之间通过配置文件和统一服务类进行衔接。配置文件负责管理数据集路径、模型路径、向量库参数和检索模式参数；服务类负责封装图像读取、特征提取、OCR 调用和结果组装逻辑。这样可以避免算法逻辑直接散落在前端接口中，也便于后续对某一模块进行替换。例如，当后续训练出独立的 \texttt{spatial\_attention} 权重后，只需要在特征提取模块中替换对应的主体增强实现，而不必改动向量数据库和前端展示逻辑。

\section{关键接口与服务组织}

为了使系统既能支持在线演示，也能支撑毕业设计中的实验复现，本文在 FastAPI 服务中设计了若干核心接口。其组织方式遵循“检索接口、管理接口、监控接口”三类职责划分。

\begin{enumerate}
  \item \textbf{检索接口。} 以 \texttt{/search} 为核心，用于接收用户上传的查询图纸，完成特征提取、向量召回和结果返回。
  \item \textbf{基础状态接口。} 包括 \texttt{/health}、\texttt{/ready}、\texttt{/live} 和 \texttt{/metrics}，分别用于健康检查、就绪检测、存活检测和监控指标输出。
  \item \textbf{图库管理接口。} 包括 \texttt{/stats}、\texttt{/api/images}、\texttt{/api/categories}、\texttt{/api/initialize}、\texttt{/api/rebuild} 等，用于图库浏览、图像管理与索引维护。
\end{enumerate}

通过上述接口组织方式，系统能够同时满足答辩演示中的交互需求和工程运行中的可维护性需求，也使系统实现与实验脚本之间形成了相对统一的接口层。

\section{索引构建实现}

早期版本中，建库流程与服务启动过程耦合较紧，且 OCR 文本抽取阶段容易造成较大的 CPU 开销。为解决这一问题，本文重新设计了独立建库流程。当前建库逻辑分为三步：
\begin{enumerate}
  \item 加载配置与模型，并显式确认视觉编码设备为 CUDA；
  \item 根据运行参数决定是否启用 YOLO 清洗和 OCR；
  \item 扫描图纸目录，按批次提取视觉特征并写入 Qdrant。
\end{enumerate}

其中，为提高 GPU 利用率，本文进一步将视觉特征提取改写为批量前向方式，不再逐张图纸单独执行 CLIP 编码。这样既提高了显卡利用效率，也降低了全量重建时的无效 CPU 占用。在正式实验前，本文已基于 \texttt{CUDA + 关闭 OCR} 的方式完成 \texttt{Dataset\_img3} 全图库重建，最终索引数量达到 13880/13880。

在工程实践中，OCR 是全量建库阶段较容易造成 CPU 压力的环节。为避免系统在首次建库时因 OCR 抽取而长时间占用 CPU，本文将视觉索引构建与 OCR 辅助评测进行解耦：视觉索引优先使用 GPU 批量完成，OCR 仅在多模态实验或需要字段辅助的候选阶段启用。该设计既保证了主检索链路的可运行性，也使不同实验模式之间的计算开销更加可控。

\section{在线检索流程实现}

系统在线检索流程如下：
\begin{enumerate}
  \item 用户上传查询图纸；
  \item 系统根据当前检索模式决定是否启用 YOLO 清洗、掩码引导聚合和结构重排序；
  \item 提取查询图纸的视觉特征与结构描述符；
  \item 在 Qdrant 中完成向量召回；
  \item 对候选结果进行结构增强与必要的 OCR 辅助重排序；
  \item 返回相似图纸路径、排序分数、结构相似度及文本辅助分析结果。
\end{enumerate}

该流程兼容论文中的双主线实验设计，即既可以运行 YOLO 清洗路线，也可以运行无清洗注意力路线。

\section{离线正式评测实现}

为了避免“图库特征来自一种模式，而查询特征来自另一种模式”导致的评测偏差，本文没有继续使用旧版仅依赖在线接口切换参数的实验方式，而是重新实现了离线模式对齐评测方案。

具体而言：
\begin{enumerate}
  \item \texttt{evaluate\_visual\_offline.py} 负责对视觉主线进行离线评测。脚本会在不同模式下分别构建图库特征缓存和查询特征，再统一计算 Recall@K、mAP@10、nDCG@10、FT、ST 和 ANMRR；
  \item \texttt{evaluate\_multimodal\_offline.py} 在视觉缓存基础上，对候选结果执行 OCR 辅助重排，从而评测 \texttt{multimodal\_text} 和 \texttt{attention\_multimodal} 两个模式；
  \item 评测脚本最终输出 JSON、CSV 和 Markdown 三种格式结果，便于论文表格回填与结果复查。
\end{enumerate}

基于该离线评测机制，本文完成了双主线 8 个模式的正式实验，并将所得结果统一回填至论文之中。

与在线检索相比，离线评测更强调公平性和可复现性。在线服务需要优先保证交互响应，而离线评测需要保证每一组实验的图库特征、查询特征和排序逻辑严格一致。因此，本文将正式指标计算全部放在线下脚本中完成，并保留原始 JSON 结果文件。这样既便于论文表格回填，也便于在答辩或复查阶段追溯每一组数值的来源。

\section{前端与交互实现}

系统前端采用 HTML 页面结合后端接口实现，支持图纸上传、相似图纸检索、图库浏览、统计信息查看和分类结果展示等功能。相较于早期简化界面，当前前端已恢复为完整展示版本，可较为完整地呈现系统功能与检索结果。用户可在页面中直接查看图库规模、检索结果、类别统计和相似样本信息，后端则负责完成特征提取、排序计算与结果返回。

在交互流程上，用户首先通过页面上传查询图纸，系统返回相似图纸缩略图、文件路径、相似度分数和相关类别信息。对于图库管理部分，页面提供类别浏览与统计入口，使用户能够了解当前图库规模和类别分布。对于系统维护部分，前端通过后端管理接口触发初始化或重建索引操作。这样的设计能够覆盖毕业设计演示中常见的三类场景：单张图纸检索、图库状态展示和系统维护操作。

\section{异常处理与运行保障}

为了提高系统运行稳定性，后端在图像读取、模型推理、向量库访问和结果返回等环节均设置了基础异常处理。当上传文件格式不符合要求、向量库尚未完成初始化或模型推理过程出现异常时，系统会返回明确的错误信息，避免前端页面长时间无响应。同时，健康检查接口可以用于判断服务是否正常启动，就绪接口可以用于判断向量库与模型是否已经完成加载。

在资源使用方面，本文将全量建库与在线查询分离。全量建库通过独立脚本执行，适合在 GPU 可用且系统空闲时运行；在线查询则只处理单张或少量图纸，避免与大规模索引构建过程相互干扰。这种实现方式能够降低系统演示过程中的不确定性，也更符合实际部署中“离线建库、在线查询”的常见工作模式。

\section{工程可运行性}

截至论文撰写完成阶段，系统已经具备较完整的工程可运行性，主要体现在：
\begin{enumerate}
  \item 已完成新数据集清理与全图库索引重建；
  \item 已支持 GPU 化建库与批量视觉编码；
  \item 已支持双主线在线检索与离线评测；
  \item 已支持 OCR 辅助检索、结构重排序和 Grad-CAM 可解释性分析；
  \item 已形成前后端联动的完整系统原型。
\end{enumerate}

上述结果表明，本文工作并未停留于算法设计层面，而是进一步形成了较为完整的系统实现与实验复现基础。

\section{本章小结}

本章从工程实现角度介绍了本文方法的落地方式。通过将数据清理、GPU 化建库、视觉特征提取、结构增强、OCR 辅助检索、离线评测和前端展示统一组织为模块化系统，本文不仅实现了工程图纸检索方法的完整运行链路，也为后续扩展独立注意力模型、优化 OCR 融合策略和部署实际系统提供了明确的实现基础。
